\section{Methodology}
\label{sec:method}

\subsection{Dataset}

Telemetry data is generated against Google's microservices-demo Kubernetes deployment with controlled fault injections from 14 distinct fault families. Each fault scenario produces logs (Loki), distributed traces (Tempo), metrics (Prometheus), and Kubernetes events; we discretize the timeline into 5-minute windows and derive 94 numeric features per window. The dataset is split chronologically and family-stratified:

\begin{table}[h]
\centering
\caption{Train/val/test splits. Scenario families do not overlap between splits, making this an out-of-distribution generalization benchmark.}
\label{tab:splits}
\begin{tabular}{lrll}
\toprule
Split & Windows & Families \\
\midrule
Train       & 2{,}796 & 14 (deploy, scheduled-job, network, etc.) \\
Validation  &   984  & 6  (checkout-restart, flapping-pod, $\ldots$) \\
Test        & 2{,}940 & 7  (ad-outage, cart-redis, currency-outage, $\ldots$) \\
\bottomrule
\end{tabular}
\end{table}

The memory corpus (Section~\ref{sec:corpus}) contains 347 V2 humanized Jira tickets. The \texttt{available\_as\_memory\_from} timestamp on each ticket enforces time-ordered visibility: a test window at time $t$ sees only tickets minted before $t$. The distractor pool contains 110 plausibly-realistic but unrelated tickets (60 \textsc{TAWOS}-derived, 25 in-architecture, 25 cross-architecture).

\paragraph{The retrievable subset.} Of 2{,}940 test windows, 317 have \texttt{gold\_label=ticket\_worthy} and a non-empty \texttt{gold\_matched\_issue\_ids} set. All retrieval metrics are scored on this subset.

\subsection{Pipelines compared}

\begin{table}[h]
\centering
\small
\caption{Pipelines compared. All memorygraph variants use the same skill chain (Section~\ref{sec:system}) with $w_{\text{num}}=0.80$; they differ in reranker choice or memory composition.}
\label{tab:pipelines}
\begin{tabular}{lll}
\toprule
Pipeline & Memory & Reranker \\
\midrule
\hgb{}                     & none                & none                        \\
\sota{}                    & V2 (347)            & MS-MARCO (off-the-shelf)    \\
\sotaft{} \textit{(ours)}  & V2 (347)            & MS-MARCO + fine-tuned       \\
\sota{}$-$logs             & V2 (logs masked)    & off-the-shelf               \\
\sota{}$-$traces           & V2 (traces masked)  & off-the-shelf               \\
\sota{}$-$k8s              & V2 (k8s masked)     & off-the-shelf               \\
\sota{}$+N\%$ distractors  & V2 + $N\%$ pool     & off-the-shelf               \\
\bottomrule
\end{tabular}
\end{table}

\subsection{Metrics}

\paragraph{Triage detection.} PR-AUC (primary), ROC-AUC, F1 at FPR=5\%, and ECE (calibration).

\paragraph{Retrieval (primary).} Hit@$K$ for $K \in \{1, 3, 5\}$. We argue Hit@$K$ is the right metric in our setting --- see Section~\ref{sec:hitk-correction}.

\paragraph{Retrieval (secondary).} MRR, Precision@5, and a capacity-normalized recall $\text{Recall@5}_{\text{norm}} = |\text{top-}5 \cap \text{gold}| / \min(5, |\text{gold}|)$.

\paragraph{Utility.} Simulated mean time-to-diagnose per resolvable incident, assuming the engineer reviews top-$K$ candidates serially at 30 seconds each and falls back to 30~minutes of manual investigation when none match.

\paragraph{Statistical envelope.} All headline numbers are paired bootstrap CIs at 95\% confidence, 1{,}000 resamples, seed 42. We also report \emph{stratified} bootstrap CIs per (pipeline, depth-bucket) cell --- the same resampling procedure applied within each stratum.

\subsection{The Hit@K methodological correction}
\label{sec:hitk-correction}

The standard $\text{Recall@}K = |\text{top-}K \cap \text{gold}| / |\text{gold}|$ definition produces $\text{Recall@}5 = 1.0$ when $|\text{gold}|=5$ and all five gold tickets are surfaced, but only $5/21 \approx 0.238$ even with perfect retrieval when $|\text{gold}|=21$. The cap is structural: at most $K$ tickets can fit in the top-$K$, so the metric is bounded by $K/|\text{gold}|$ whenever $|\text{gold}|>K$.

In our setting, $|\text{gold}|$ varies from 1 to 30+ across the test split because incidents with rich deployment history have many compatible prior tickets. This variation makes Recall@$K$ \emph{incomparable across depth buckets}: the depth-stratified curve falsely drops from 0.155 at depth 1--2 to 0.044 at depth 21+ in our anchor experiment, simply because the denominator grows faster than $K$.

We argue Hit@$K = \mathbb{1}[\text{top-}K \cap \text{gold} \neq \emptyset]$ is the right metric for ``did the engineer find a relevant ticket.'' Hit@$K$ has a clean operational semantics: it answers ``with probability $p$, will a user scanning top-$K$ candidates see at least one useful match?'' This is what matters for the engineer-time-saved framing, and it produces strictly monotone depth-scaling curves in our data (Section~\ref{sec:results}). We report Hit@$K$ as primary and the capacity-normalized $\text{Recall@}K_{\text{norm}}$ as secondary for completeness.

\subsection{Reproducibility}

Every pipeline fit is wrapped in an \texttt{open\_training\_run} context manager that creates a directory \texttt{data/derived/global/<id>/training\_runs/\,<pipeline>\_\_<UTC>\_\_<sha8>/} and writes:

\begin{itemize}
    \item \texttt{config.json} --- pipeline class, hyperparameters, target FPR, paths, Python version, git SHA.
    \item \texttt{metrics.json} --- threshold, fit/predict timing, sample count.
    \item \texttt{predictions.jsonl} --- per-window predictions matching the \texttt{PipelinePrediction} schema.
    \item \texttt{train.log} --- human-readable lifecycle log in UTC.
\end{itemize}

Every reported number in this paper can be traced to a specific git commit through this directory's \texttt{config.json}. The paper's anchor figures come from commit \texttt{f4684b3}; Phases B--D append commits \texttt{742e9e5}, \texttt{32297cb}, \texttt{d970e5f}.
